\chapter{Hilbert空間}

\paragraph{本章の目標}
前章では，関数の大きさを測るノルムを導入し，
$L^p$ 空間や $C(X)$ がBanach空間になることを見た．
その中でも
\[
p=2
\]
の場合は特別である．
このときノルムは内積から作られ，
長さだけでなく角度や直交も扱えるようになる．

本章では
\begin{enumerate}
    \item 内積空間・Hilbert空間の定義
    \item Cauchy--Schwarz の不等式と内積から誘導されるノルム
    \item $L^2(X,m)$ がHilbert空間になること
    \item 正規直交系，Fourier係数，Bessel の不等式，Parseval の等式
    \item Fourier関数系という典型例
\end{enumerate}
を説明する．

\section{動機づけ}

ノルムだけでは，ベクトルの長さは分かっても，
2つのベクトルがどれくらい「直交しているか」は分からない．
しかし実際には
\[
\text{最小二乗法},\qquad
\text{Fourier級数},\qquad
\text{エネルギー評価}
\]
のように，直交が本質的な役割を果たす場面が多い．

例えば $L^2([0,1])$ の関数 $f$ を
\[
1,\ \cos(2\pi x),\ \sin(2\pi x),\dots
\]
のような基本関数で展開したいとき，
係数は
\[
\int_0^1 f(x)\cos(2\pi nx)\,dx,
\qquad
\int_0^1 f(x)\sin(2\pi nx)\,dx
\]
で与えられる．
これらは単なる積分ではなく，
関数 $f$ と基本関数との「内積」になっている．

この考え方を抽象化したのがHilbert空間である．

\section{内積空間とHilbert空間}

\begin{definition}[内積]
実ベクトル空間 $H$ 上の写像
\[
\langle \cdot,\cdot\rangle:H\times H\to\RR
\]
が内積であるとは，任意の $x,y,z\in H$ と任意の $a\in\RR$ に対して
次を満たすことをいう．
\begin{enumerate}
    \item
    \[
    \langle x+ y,z\rangle=\langle x,z\rangle+\langle y,z\rangle
    \]
    \item
    \[
    \langle ax,y\rangle=a\langle x,y\rangle
    \]
    \item
    \[
    \langle x,y\rangle=\langle y,x\rangle
    \]
    \item
    \[
    \langle x,x\rangle\ge 0
    \]
    かつ
    \[
    \langle x,x\rangle=0 \iff x=0
    \]
\end{enumerate}
\end{definition}

\begin{definition}[内積空間]
内積を備えた実ベクトル空間を内積空間という．
\end{definition}

\begin{example}[\texorpdfstring{$\RR^d$}{Rd} の標準内積]
$x=(x_1,\dots,x_d)$，$y=(y_1,\dots,y_d)\in\RR^d$ に対して
\[
\langle x,y\rangle:=\sum_{k=1}^d x_k y_k
\]
と定めると，これは $\RR^d$ 上の内積である．
\end{example}

\begin{theorem}[Cauchy--Schwarzの不等式]
内積空間 $H$ の任意の $x,y\in H$ に対して
\[
|\langle x,y\rangle|
\le
\sqrt{\langle x,x\rangle}\sqrt{\langle y,y\rangle}
\]
が成り立つ．
\end{theorem}
\begin{proof}
$y=0$ のときは自明である．
以下では $y\ne 0$ とする．
任意の $t\in\RR$ に対して
\[
\langle x-ty,x-ty\rangle\ge 0
\]
である．
これを展開すると
\[
\langle x,x\rangle-2t\langle x,y\rangle+t^2\langle y,y\rangle\ge 0
\]
を得る．
左辺は $t$ の2次式であり，すべての $t$ に対して非負だから判別式は $0$ 以下である．
したがって
\[
4\langle x,y\rangle^2-4\langle x,x\rangle\langle y,y\rangle\le 0
\]
すなわち
\[
\langle x,y\rangle^2\le \langle x,x\rangle\langle y,y\rangle
\]
が成り立つ．
両辺は非負だから平方根をとって結論を得る．
\end{proof}

\begin{corollary}[内積から誘導されるノルム]
内積空間 $H$ に対して
\[
\|x\|:=\sqrt{\langle x,x\rangle}
\]
とおくと，これは $H$ 上のノルムである．
\end{corollary}
\begin{proof}
まず
\[
\|x\|=0
\iff
\langle x,x\rangle=0
\iff
x=0
\]
である．

次に
\[
\|ax\|^2=\langle ax,ax\rangle=a^2\langle x,x\rangle=|a|^2\|x\|^2
\]
だから
\[
\|ax\|=|a|\,\|x\|
\]
である．

最後に
\begin{align*}
\|x+y\|^2
&=
\langle x+y,x+y\rangle \\
&=
\|x\|^2+2\langle x,y\rangle+\|y\|^2 \\
&\le
\|x\|^2+2\|x\|\|y\|+\|y\|^2 \\
&=
(\|x\|+\|y\|)^2
\end{align*}
である．
両辺は非負だから平方根をとって
\[
\|x+y\|\le \|x\|+\|y\|
\]
を得る．
\end{proof}

\begin{proposition}[Pythagorasの定理]
内積空間 $H$ の $x,y\in H$ が
\[
\langle x,y\rangle=0
\]
を満たすとする．
このとき
\[
\|x+y\|^2=\|x\|^2+\|y\|^2
\]
が成り立つ．
\end{proposition}
\begin{proof}
\[
\|x+y\|^2
=
\langle x+y,x+y\rangle
=
\langle x,x\rangle+\langle x,y\rangle+\langle y,x\rangle+\langle y,y\rangle
=
\|x\|^2+\|y\|^2
\]
である．
\end{proof}

\begin{definition}[Hilbert空間]
内積空間 $H$ が，
上の内積から誘導されるノルムに関して完備であるとき，
$H$ をHilbert空間という．
\end{definition}

\begin{example}
$\RR^d$ は標準内積に関してHilbert空間である．
実際，誘導されるノルムは
\[
\|x\|=\left(\sum_{k=1}^d |x_k|^2\right)^{1/2}
\]
であり，前章で $\RR^d$ はこのノルムに関して完備であることを示した．
\end{example}

\section{\texorpdfstring{$L^2(X,m)$}{L2(X,m)} はHilbert空間}

以後，
\[
(X,\calF,m)
\]
を測度空間とする．

\begin{definition}[\texorpdfstring{$L^2$}{L2} の内積]
$f,g\in L^2(X,m)$ に対して
\[
\langle f,g\rangle:=\int_X f g\,dm
\]
と定める．
\end{definition}

\begin{proposition}
上の式は $L^2(X,m)$ 上のwell-definedな内積である．
\end{proposition}
\begin{proof}
まず $f,g\in L^2(X,m)$ なら，
前章の Hölder の不等式を $p=q=2$ に適用して
\[
\int_X |fg|\,dm\le \|f\|_2\|g\|_2<\infty
\]
を得る．
したがって積分は有限である．

また $f=f'$ a.e., $g=g'$ a.e.\ なら
\[
fg=f'g' \quad \text{a.e.}
\]
だから，第8章の a.e.\ 不変性より
\[
\int_X fg\,dm=\int_X f'g'\,dm
\]
である．
よってこの式は同値類の取り方によらない．

線形性と対称性は積分の線形性から直ちに従う．
さらに
\[
\langle f,f\rangle=\int_X f^2\,dm\ge 0
\]
である．
もし
\[
\langle f,f\rangle=0
\]
なら
\[
\int_X f^2\,dm=0
\]
であり，第9章の「積分が $0$ であることの判定」より
\[
f^2=0 \quad \text{a.e.}
\]
したがって $f=0$ a.e.\ である．
すなわち $L^2$ の元として $f=0$ である．
以上より，これは内積である．
\end{proof}

\begin{theorem}
\[
L^2(X,m)
\]
は上の内積に関してHilbert空間である．
\end{theorem}
\begin{proof}
直前の命題より，$L^2(X,m)$ は内積空間である．
この内積が誘導するノルムは
\[
\|f\|=\sqrt{\langle f,f\rangle}
=
\left(\int_X |f|^2\,dm\right)^{1/2}
=
\|f\|_2
\]
である．
前章の Riesz--Fischer の定理より，
$L^2(X,m)$ はこのノルムに関して完備である．
したがって $L^2(X,m)$ はHilbert空間である．
\end{proof}

\begin{example}
$X=[0,1]$ とし $m$ をLebesgue測度とすると，
\[
L^2([0,1],m)
\]
はHilbert空間である．
ここでの内積は
\[
\langle f,g\rangle=\int_0^1 f(x)g(x)\,dx
\]
である．
確率論では，確率空間上の $L^2$ は「二乗平均が有限な確率変数全体」の空間になる．
\end{example}

\section{正規直交系と直交射影}

\begin{definition}[直交と正規直交]
内積空間 $H$ の $x,y\in H$ に対して
\[
\langle x,y\rangle=0
\]
のとき，$x$ と $y$ は直交するといい，
\[
x\perp y
\]
と書く．

また $\|x\|=1$ を満たすベクトルを正規化されたベクトルという．
互いに直交し，かつ各ベクトルが長さ $1$ をもつベクトル族を正規直交族という．
\end{definition}

\begin{definition}[正規直交系とCONS]
Hilbert空間 $H$ の列
\[
\{e_n\}_{n=1}^{\infty}
\]
が正規直交系であるとは，
任意の $m,n\in\NN$ に対して
\[
\langle e_n,e_m\rangle=
\begin{cases}
1 & (n=m),\\
0 & (n\ne m)
\end{cases}
\]
が成り立つことをいう．

さらに
\[
\operatorname{span}\{e_1,e_2,\dots\}
\]
の閉包が $H$ 全体に一致するとき，
この正規直交系を完全正規直交系
（complete orthonormal system, CONS）という．
\end{definition}

\begin{proposition}
正規直交族は線形独立である．
\end{proposition}
\begin{proof}
\[
\sum_{k=1}^N c_k e_k=0
\]
とする．
両辺と $e_j$ との内積をとると
\[
0
=
\left\langle \sum_{k=1}^N c_k e_k,e_j\right\rangle
=
\sum_{k=1}^N c_k \langle e_k,e_j\rangle
=
c_j
\]
である．
$j=1,\dots,N$ は任意だから，すべての係数 $c_j$ は $0$ である．
\end{proof}

\begin{definition}[Fourier係数と部分和]
$\{e_n\}_{n=1}^{\infty}$ をHilbert空間 $H$ の正規直交系とする．
$x\in H$ に対して
\[
\langle x,e_n\rangle
\]
を $x$ の $e_n$ に関するFourier係数という．

また
\[
P_Nx:=\sum_{n=1}^N \langle x,e_n\rangle e_n
\]
を $x$ の第 $N$ 部分和という．
\end{definition}

\begin{theorem}[直交射影の基本性質]
$\{e_n\}$ を正規直交系とし，
\[
M_N:=\operatorname{span}\{e_1,\dots,e_N\}
\]
とおく．
このとき任意の $x\in H$ に対して次が成り立つ．
\begin{enumerate}
    \item
    \[
    x-P_Nx \perp e_j
    \qquad (1\le j\le N)
    \]
    \item 任意の $y\in M_N$ に対して
    \[
    \|x-y\|^2
    =
    \|x-P_Nx\|^2
    +
    \|P_Nx-y\|^2
    \]
    が成り立つ．
    特に
    \[
    \|x-P_Nx\|\le \|x-y\|
    \]
    である．
\end{enumerate}
\end{theorem}
\begin{proof}
\textbf{(1)}\quad
$1\le j\le N$ に対して
\begin{align*}
\langle x-P_Nx,e_j\rangle
&=
\langle x,e_j\rangle
-
\left\langle \sum_{n=1}^N \langle x,e_n\rangle e_n,e_j\right\rangle \\
&=
\langle x,e_j\rangle
-
\sum_{n=1}^N \langle x,e_n\rangle \langle e_n,e_j\rangle \\
&=
\langle x,e_j\rangle-\langle x,e_j\rangle
=0.
\end{align*}

\textbf{(2)}\quad
$y\in M_N$ を任意に取る．
すると
\[
P_Nx-y\in M_N
\]
であるから，(1) より
\[
x-P_Nx \perp P_Nx-y
\]
が成り立つ．
さらに
\[
x-y=(x-P_Nx)+(P_Nx-y)
\]
だから，Pythagoras の定理より
\[
\|x-y\|^2
=
\|x-P_Nx\|^2+\|P_Nx-y\|^2
\]
を得る．
特に第2項は非負だから
\[
\|x-P_Nx\|\le \|x-y\|
\]
である．
\end{proof}

\begin{corollary}[Besselの不等式]
$\{e_n\}_{n=1}^{\infty}$ をHilbert空間 $H$ の正規直交系とする．
このとき任意の $x\in H$ と任意の $N\in\NN$ に対して
\[
\sum_{n=1}^N |\langle x,e_n\rangle|^2\le \|x\|^2
\]
が成り立つ．
特に級数
\[
\sum_{n=1}^{\infty} |\langle x,e_n\rangle|^2
\]
は収束する．
\end{corollary}
\begin{proof}
直交射影の基本性質を $y=0$ に適用すると
\[
\|x\|^2=\|x-P_Nx\|^2+\|P_Nx\|^2
\]
を得る．
一方，$\{e_1,\dots,e_N\}$ は正規直交族だから
\[
\|P_Nx\|^2
=
\left\|
\sum_{n=1}^N \langle x,e_n\rangle e_n
\right\|^2
=
\sum_{n=1}^N |\langle x,e_n\rangle|^2
\]
である．
したがって
\[
\sum_{n=1}^N |\langle x,e_n\rangle|^2
=
\|P_Nx\|^2
\le
\|x\|^2
\]
である．
左辺は $N$ に関して単調増加で上に有界だから収束する．
\end{proof}

\begin{theorem}[CONSの特徴づけ]
$\{e_n\}_{n=1}^{\infty}$ をHilbert空間 $H$ の正規直交系とする．
このとき次は同値である．
\begin{enumerate}
    \item $\{e_n\}$ は CONS である
    \item 任意の $x\in H$ に対して
    \[
    \|x-P_Nx\|\to 0
    \]
    が成り立つ
    \item 任意の $x\in H$ に対して
    \[
    \|x\|^2=\sum_{n=1}^{\infty}|\langle x,e_n\rangle|^2
    \]
    が成り立つ
\end{enumerate}
\end{theorem}
\begin{proof}
\textbf{(1) $\Rightarrow$ (2)}\quad
$x\in H$ と $\eps>0$ を任意に取る．
CONS であることから，ある有限和
\[
y=\sum_{n=1}^M a_n e_n
\]
が存在して
\[
\|x-y\|<\eps
\]
である．
$N\ge M$ とすれば $y\in M_N$ であるから，直交射影の基本性質より
\[
\|x-P_Nx\|\le \|x-y\|<\eps.
\]
よって
\[
\|x-P_Nx\|\to 0
\]
である．

\textbf{(2) $\Rightarrow$ (3)}\quad
Bessel の不等式の証明で得た等式
\[
\|x\|^2=\|x-P_Nx\|^2+\sum_{n=1}^N |\langle x,e_n\rangle|^2
\]
で $N\to\infty$ とすれば，
\[
\|x\|^2=\sum_{n=1}^{\infty}|\langle x,e_n\rangle|^2
\]
を得る．

\textbf{(3) $\Rightarrow$ (2)}\quad
同じ等式に (3) を代入すると
\[
\|x-P_Nx\|^2
=
\sum_{n=N+1}^{\infty}|\langle x,e_n\rangle|^2
\to 0
\]
である．
したがって
\[
\|x-P_Nx\|\to 0.
\]

\textbf{(2) $\Rightarrow$ (1)}\quad
各 $P_Nx$ は
\[
\operatorname{span}\{e_1,e_2,\dots\}
\]
に属する．
(2) より $P_Nx\to x$ だから，
$x$ はその閉包に属する．
$x\in H$ は任意だから，この閉包は $H$ 全体に一致する．
したがって $\{e_n\}$ は CONS である．
\end{proof}

\begin{example}[\texorpdfstring{$\RR^d$}{Rd} の標準基底]
$\RR^d$ の標準基底
\[
e_1=(1,0,\dots,0),\ \dots,\ e_d=(0,\dots,0,1)
\]
は正規直交族であり，
その線形包は $\RR^d$ 全体だから CONS である．
このとき
\[
x=\sum_{k=1}^d \langle x,e_k\rangle e_k
\]
であり，
Parseval の等式は
\[
\|x\|^2=\sum_{k=1}^d |x_k|^2
\]
に他ならない．
\end{example}

\section{Fourier関数系}

\begin{definition}[実Fourier関数系]
$L^2([0,1],m)$ において
\[
\varphi_0(x):=1,
\qquad
\varphi_n^c(x):=\sqrt{2}\cos(2\pi n x),
\qquad
\varphi_n^s(x):=\sqrt{2}\sin(2\pi n x)
\quad (n\ge 1)
\]
を考える．
これらをまとめて実Fourier関数系という．
\end{definition}

\begin{proposition}
実Fourier関数系は $L^2([0,1],m)$ の正規直交系である．
\end{proposition}
\begin{proof}
まず
\[
\int_0^1 1^2\,dx=1
\]
だから $\varphi_0$ は長さ $1$ をもつ．

次に $n\ge 1$ とすると
\[
\int_0^1 \cos^2(2\pi n x)\,dx
=
\frac12\int_0^1 \bigl(1+\cos(4\pi n x)\bigr)\,dx
=
\frac12
\]
であるから
\[
\|\varphi_n^c\|_2^2=2\cdot \frac12=1.
\]
同様に
\[
\int_0^1 \sin^2(2\pi n x)\,dx
=
\frac12\int_0^1 \bigl(1-\cos(4\pi n x)\bigr)\,dx
=
\frac12
\]
より
\[
\|\varphi_n^s\|_2^2=1.
\]

また任意の非零整数 $k$ に対して
\[
\int_0^1 \cos(2\pi kx)\,dx=0,
\qquad
\int_0^1 \sin(2\pi kx)\,dx=0
\]
である．
したがって
\[
\langle \varphi_0,\varphi_n^c\rangle
=
\sqrt{2}\int_0^1 \cos(2\pi n x)\,dx
=0,
\]
\[
\langle \varphi_0,\varphi_n^s\rangle
=
\sqrt{2}\int_0^1 \sin(2\pi n x)\,dx
=0.
\]

さらに $m\ne n$ に対して積和公式
\[
\cos \alpha \cos \beta
=
\frac12\bigl(\cos(\alpha-\beta)+\cos(\alpha+\beta)\bigr)
\]
を用いると
\begin{align*}
\langle \varphi_n^c,\varphi_m^c\rangle
&=
2\int_0^1 \cos(2\pi n x)\cos(2\pi m x)\,dx \\
&=
\int_0^1 \cos(2\pi(n-m)x)\,dx
+
\int_0^1 \cos(2\pi(n+m)x)\,dx \\
&=0.
\end{align*}
$n=m$ のときはすでに長さ $1$ を確かめたので
\[
\langle \varphi_n^c,\varphi_n^c\rangle=1
\]
である．

同様に
\[
\sin \alpha \sin \beta
=
\frac12\bigl(\cos(\alpha-\beta)-\cos(\alpha+\beta)\bigr)
\]
より
\[
\langle \varphi_n^s,\varphi_m^s\rangle
=
\delta_{nm}
\]
を得る．

最後に
\[
\sin \alpha \cos \beta
=
\frac12\bigl(\sin(\alpha+\beta)+\sin(\alpha-\beta)\bigr)
\]
より
\[
\langle \varphi_n^s,\varphi_m^c\rangle=0
\]
である．
したがって実Fourier関数系は正規直交系である．
\end{proof}

\begin{remark}
実Fourier関数系に対するFourier係数は
\[
\langle f,\varphi_0\rangle=\int_0^1 f(x)\,dx,
\]
\[
\langle f,\varphi_n^c\rangle
=
\sqrt{2}\int_0^1 f(x)\cos(2\pi n x)\,dx,
\qquad
\langle f,\varphi_n^s\rangle
=
\sqrt{2}\int_0^1 f(x)\sin(2\pi n x)\,dx
\]
で与えられる．
つまり Fourier 係数は「関数と基本波との内積」である．

古典的な Fourier 解析では，
この正規直交系が実際に CONS であることを示す．
その結果，$L^2([0,1])$ の関数をFourier級数で展開できる．
\end{remark}

\section{解釈とまとめ}

Banach空間では長さだけを見ていたが，
Hilbert空間では内積によって
\[
\text{長さ},\qquad \text{角度},\qquad \text{直交}
\]
が同時に扱えるようになった．

特に重要なのは，
\[
L^2(X,m)
\]
がHilbert空間になることである．
ここでは関数の直交性が積分
\[
\int_X f g\,dm
\]
で表されるため，
積分論と線形代数が直接結びつく．

さらに正規直交系を用いると，
任意のベクトル $x$ を
\[
\langle x,e_n\rangle
\]
という係数で解析できる．
Bessel の不等式と Parseval の等式は，
この展開がエネルギー保存
\[
\|x\|^2
\]
と整合的であることを示している．
次章では，こうした空間のあいだを結ぶ写像として
線形作用素を導入する．
